% Axioms of GPA, in addition to those of GAA: the inequality generator is a
% lax (co)monoid morphism (P1-P4), interacts with scalars by sign (P5, P6),
% and is antisymmetric, spider-like and directed.
\begin{center}\begin{tabular}{r c c c}
  (P1) & \begin{tikzpicture}[axpic]
      \pic (g1) at (0.500,0.500) {ge};
      \pic (g2) at (1.750,0.500) {copy};
      \draw (g1-out-0) to[out=0,in=180] (g2-in-0);
  \end{tikzpicture}
  & \;\begin{tikzpicture}[axpic]\node{$\subseteq$};\end{tikzpicture}\;
  & \begin{tikzpicture}[axpic]
      \pic (g1) at (0.500,0.500) {copy};
      \pic (g2) at (1.750,1.000) {ge};
      \pic (g3) at (1.750,0.000) {ge};
      \draw (g1-out-0) to[out=0,in=180] (g2-in-0);
      \draw (g1-out-1) to[out=0,in=180] (g3-in-0);
  \end{tikzpicture} \\[1mm]
  \diagrow{(P2)}{
      \pic (g1) at (0.500,0.500) {multiplication};
      \pic (g2) at (1.750,0.500) {ge};
      \draw (g1-out-0) to[out=0,in=180] (g2-in-0);
  }{
      \pic (g1) at (0.500,1.000) {ge};
      \pic (g2) at (0.500,0.000) {ge};
      \pic (g3) at (1.750,0.500) {multiplication};
      \draw (g1-out-0) to[out=0,in=180] (g3-in-0);
      \draw (g2-out-0) to[out=0,in=180] (g3-in-1);
  }
  \diagrow{(P3)}{
      \pic (g1) at (0.500,0.500) {ge};
      \pic (g2) at (1.750,0.500) {discard};
      \draw (g1-out-0) -- (g2-in-0);
  }{
      \pic (g1) at (0.500,0.500) {discard};
      \draw (-0.250,0.500) -- (g1-in-0);
  }
  (P4) & \begin{tikzpicture}[axpic]
      \pic (g1) at (0.500,0.500) {zero};
      \draw (g1-out-0) -- (1.250,0.500);
  \end{tikzpicture}
  & \;\begin{tikzpicture}[axpic]\node{$\subseteq$};\end{tikzpicture}\;
  & \begin{tikzpicture}[axpic]
      \pic (g1) at (0.500,0.500) {zero};
      \pic (g2) at (1.750,0.500) {ge};
      \draw (g1-out-0) to[out=0,in=180] (g2-in-0);
  \end{tikzpicture} \\[1mm]
\end{tabular}\end{center}

\begin{center}\begin{tabular}{r c c c}
  \diagrow{(P5) \ $k>0$}{
      \pic[val=k] (g1) at (0.500,0.500) {scalar};
      \pic (g2) at (1.750,0.500) {ge};
      \draw (g1-out-0) to[out=0,in=180] (g2-in-0);
  }{
      \pic (g1) at (0.500,0.500) {ge};
      \pic[val=k] (g2) at (1.750,0.500) {scalar};
      \draw (g1-out-0) to[out=0,in=180] (g2-in-0);
  }
  \diagrow{(P6) \ $k<0$}{
      \pic[val=k] (g1) at (0.500,0.500) {scalar};
      \pic (g2) at (1.750,0.500) {ge};
      \draw (g1-out-0) to[out=0,in=180] (g2-in-0);
  }{
      \pic (g1) at (0.500,0.500) {coge};
      \pic[val=k] (g2) at (1.750,0.500) {scalar};
      \draw (g1-out-0) to[out=0,in=180] (g2-in-0);
  }
\end{tabular}\end{center}

\begin{center}\begin{tabular}{r c c c}
  (antisym) & \begin{tikzpicture}[axpic]
      \pic (g1) at (0.500,0.500) {copy};
      \pic (g2) at (1.750,1.000) {ge};
      \pic (g3) at (1.750,0.000) {coge};
      \pic (g4) at (3.000,0.500) {cocopy};
      \draw (g1-out-0) to[out=0,in=180] (g2-in-0);
      \draw (g1-out-1) to[out=0,in=180] (g3-in-0);
      \draw (g2-out-0) to[out=0,in=180] (g4-in-0);
      \draw (g3-out-0) to[out=0,in=180] (g4-in-1);
  \end{tikzpicture}
  & \;\begin{tikzpicture}[axpic]\node{$\subseteq$};\end{tikzpicture}\;
  & \begin{tikzpicture}[axpic]
      \draw (0.000,0.500) -- (1.250,0.500);
  \end{tikzpicture} \\[2mm]
  \diagrow{(spider)}{
      \pic (g1) at (0.500,1.000) {ge};
      \pic (g2) at (0.500,0.000) {ge};
      \pic (c1) at (1.750,0.500) {cocopy};
      \pic (c2) at (3.000,0.500) {copy};
      \pic (h1) at (4.250,1.000) {ge};
      \pic (h2) at (4.250,0.000) {ge};
      \draw (g1-out-0) to[out=0,in=180] (c1-in-0);
      \draw (g2-out-0) to[out=0,in=180] (c1-in-1);
      \draw (c1-out-0) -- (c2-in-0);
      \draw (c2-out-0) to[out=0,in=180] (h1-in-0);
      \draw (c2-out-1) to[out=0,in=180] (h2-in-0);
  }{
      \pic (c1) at (0.500,1.500) {copy};
      \pic (c2) at (0.500,-0.500) {copy};
      \pic (g11) at (2.000,2.000) {ge};
      \pic (g12) at (2.000,1.000) {ge};
      \pic (g21) at (2.000,0.000) {ge};
      \pic (g22) at (2.000,-1.000) {ge};
      \pic (d1) at (3.750,1.500) {cocopy};
      \pic (d2) at (3.750,-0.500) {cocopy};
      \draw (c1-out-0) to[out=0,in=180] (g11-in-0);
      \draw (c1-out-1) to[out=0,in=180] (g12-in-0);
      \draw (c2-out-0) to[out=0,in=180] (g21-in-0);
      \draw (c2-out-1) to[out=0,in=180] (g22-in-0);
      \draw (g11-out-0) to[out=0,in=180] (d1-in-0);
      \draw (g21-out-0) to[out=0,in=180] (d1-in-1);
      \draw (g12-out-0) to[out=0,in=180] (d2-in-0);
      \draw (g22-out-0) to[out=0,in=180] (d2-in-1);
  }
  \diagrow{(direction)}{
      \pic (g1) at (0.500,0.500) {coge};
      \pic (g2) at (1.500,0.500) {ge};
      \draw (g1-out-0) -- (g2-in-0);
  }{
      \pic (g1) at (0.500,0.500) {discard};
      \pic (g2) at (1.750,0.500) {codiscard};
      \draw (-0.250,0.500) -- (g1-in-0);
      \draw (g2-out-0) -- (2.500,0.500);
  }
\end{tabular}\end{center}
